% Mathematical TeX extracted from the cited web pages, reformatted by the assistant. Not the native full chapter.

\begin{lemma}[10.115.1, Tag 051M]
Let $n\in\mathbf N$. Let $N$ be a finite nonempty set of multi-indices $\nu=(\nu_1,\ldots,\nu_n)$. Given $e=(e_1,\ldots,e_n)$ we set $e\cdot\nu=\sum e_i\nu_i$. Then for $e_1\gg e_2\gg\ldots\gg e_{n-1}\gg e_n$ we have: If $\nu,\nu'\in N$ then
\[
(e\cdot\nu=e\cdot\nu')\Longleftrightarrow(\nu=\nu').
\]
\end{lemma}
\begin{proof}
Say $N=\{\nu_j\}$ with $\nu_j=(\nu_{j1},\ldots,\nu_{jn})$. Let $A_i=\max_j\nu_{ji}-\min_j\nu_{ji}$. If for each $i$ we have $e_{i-1}>A_ie_i+A_{i+1}e_{i+1}+\ldots+A_ne_n$ then the lemma holds. For suppose that $e\cdot(\nu-\nu')=0$. Then for $n\geq2$,
\[
e_1(\nu_1-\nu'_1)=\sum_{i=2}^n e_i(\nu'_i-\nu_i).
\]
We may assume that $(\nu_1-\nu'_1)\geq0$. If $(\nu_1-\nu'_1)>0$, then
\begin{align*}
e_1(\nu_1-\nu'_1)&\geq e_1>A_2e_2+\ldots+A_ne_n\\
&\geq\sum_{i=2}^n e_i|\nu'_i-\nu_i|\geq\sum_{i=2}^n e_i(\nu'_i-\nu_i).
\end{align*}
This contradiction implies that $\nu'_1=\nu_1$. By induction, $\nu'_i=\nu_i$ for $2\leq i\leq n$.
\end{proof}
\begin{lemma}[10.115.2, Tag 051N]
Let $R$ be a ring. Let $g\in R[x_1,\ldots,x_n]$ be an element which is nonconstant, i.e., $g\notin R$. For $e_1\gg e_2\gg\ldots\gg e_{n-1}\gg e_n=1$ the polynomial
\[
g(x_1+x_n^{e_1},x_2+x_n^{e_2},\ldots,x_{n-1}+x_n^{e_{n-1}},x_n)
=ax_n^d+\text{lower order terms in }x_n
\]
where $d>0$ and $a\in R$ is one of the nonzero coefficients of $g$.
\end{lemma}
\begin{proof}
Write $g=\sum_{\nu\in N}a_\nu x^\nu$ with $a_\nu\in R$ not zero. Here $N$ is a finite set of multi-indices as in Lemma 10.115.1 and $x^\nu=x_1^{\nu_1}\ldots x_n^{\nu_n}$. Note that the leading term in
\[
(x_1+x_n^{e_1})^{\nu_1}\ldots(x_{n-1}+x_n^{e_{n-1}})^{\nu_{n-1}}x_n^{\nu_n}
\quad\text{is}\quad x_n^{e_1\nu_1+\ldots+e_{n-1}\nu_{n-1}+\nu_n}.
\]
Hence the lemma follows from Lemma 10.115.1 which guarantees that there is exactly one nonzero term $a_\nu x^\nu$ of $g$ which gives rise to the leading term of $g(x_1+x_n^{e_1},x_2+x_n^{e_2},\ldots,x_{n-1}+x_n^{e_{n-1}},x_n)$, i.e., $a=a_\nu$ for the unique $\nu\in N$ such that $e\cdot\nu$ is maximal.
\end{proof}
\begin{lemma}[10.115.3, Tag 00OX]
Let $k$ be a field. Let $S=k[x_1,\ldots,x_n]/I$ for some proper ideal $I$. If $I\neq0$, then there exist $y_1,\ldots,y_{n-1}\in k[x_1,\ldots,x_n]$ such that $S$ is finite over $k[y_1,\ldots,y_{n-1}]$. Moreover we may choose $y_i$ to be in the $\mathbf Z$-subalgebra of $k[x_1,\ldots,x_n]$ generated by $x_1,\ldots,x_n$.
\end{lemma}
\begin{proof}
Pick $f\in I$, $f\neq0$. It suffices to show the lemma for $k[x_1,\ldots,x_n]/(f)$ since $S$ is a quotient of that ring. We will take $y_i=x_i-x_n^{e_i}$, $i=1,\ldots,n-1$ for suitable integers $e_i$. When does this work? It suffices to show that $\overline{x_n}\in k[x_1,\ldots,x_n]/(f)$ is integral over the ring $k[y_1,\ldots,y_{n-1}]$. The equation for $\overline{x_n}$ over this ring is
\[
f(y_1+x_n^{e_1},\ldots,y_{n-1}+x_n^{e_{n-1}},x_n)=0.
\]
Hence we are done if we can show there exists integers $e_i$ such that the leading coefficient with respect to $x_n$ of the equation above is a nonzero element of $k$. This can be achieved for example by choosing $e_1\gg e_2\gg\ldots\gg e_{n-1}$, see Lemma 10.115.2.
\end{proof}
\begin{lemma}[10.115.4, Noether normalization, Tag 00OY]
Let $k$ be a field. Let $S=k[x_1,\ldots,x_n]/I$ for some ideal $I$. If $I\neq(1)$, there exist $r\geq0$, and $y_1,\ldots,y_r\in k[x_1,\ldots,x_n]$ such that (a) the map $k[y_1,\ldots,y_r]\to S$ is injective where the source is the polynomial ring on $y_1,\ldots,y_r$, and (b) the map $k[y_1,\ldots,y_r]\to S$ is finite. In this case the integer $r$ is the dimension of $S$.

Moreover we may choose $y_i$ to be in the $\mathbf Z$-subalgebra of $k[x_1,\ldots,x_n]$ generated by $x_1,\ldots,x_n$.
\end{lemma}
\begin{proof}
By induction on $n$, with $n=0$ being trivial. If $I=0$, then take $r=n$ and $y_i=x_i$. If $I\neq0$, then choose $y_1,\ldots,y_{n-1}$ as in Lemma 10.115.3. Let $S'\subset S$ be the subring generated by the images of the $y_i$. By induction we can choose $r$ and $z_1,\ldots,z_r\in k[y_1,\ldots,y_{n-1}]$ such that (a), (b) hold for $k[z_1,\ldots,z_r]\to S'$. Since $S'\to S$ is injective and finite we see (a), (b) hold for $k[z_1,\ldots,z_r]\to S$. The last assertion follows from Lemma 10.112.4.
\end{proof}

